\subsection{No Adaptive Scheme}\label{sec:noadapt}
Suppose we have a set of $N$ sample points $X = [x_1, x_2, ..., x_N]$, independently drawn from a Gaussian mixture $f_w(x) = \sum_{i = 1}^p b_i(x)w_i(x)$, which is a weighted combination of $p$ fixed Gaussian basis functions. For the first step, let's define the objective as to estimate the weights $w_{(p \times 1)} = [w_1, w_2, ..., w_p]^T$ in the $f_w$, knowing $p$ fixed basis functions $b_{(p\times 1)} = [b_1(\cdot), b_2(\cdot), ..., b_p(\cdot)]^T$.\par
The idea is to assume a set of (kernel) functions $k_{(m \times 1)} = [k_1(\cdot), k_2(\cdot),..., k_m(\cdot)]^T$ and define the new random variable $y$ as:
\begin{equation}
    y_{(m \times 1)} = 1/N \sum_{i=1}^N k(x_i)
\end{equation}
because $y$ is sum of some random variables with the same distribution, it ($y$) has a Gaussian distribution (Central limit theorem). Our new aim is then to estimate the distribution of $y$ using the maximum likelihood method. The maximum likelihood objective is to:
\begin{equation} \label{eq:prop_obj1}
    min_{w} 1/2(E(y) - y)^TC^{-1}(E(y) - y)
\end{equation}
where
\begin{align}
    E(y)_{(m \times 1)} = \int k(x)f_w(x) dx \\
    = \left[\int k_{(m \times 1)}(x)b^T_{(1 \times p)}(x)dx \right]w_{{p \times1}} \\
    =Dw
\end{align}
and $C^{-1}_{(m \times M)}$ is the covariance matrix of $y$ which we compute it empirically. For computing $C^{-1}$, we draw $N$ samples from the distribution for $T$ times and compute $T$ random variables $y$ and then compute the covariance matrix of $T$ computed $ys$. \par
 We also add a sparsity term to the objective \eqref{eq:prop_obj1} as :
 \begin{align}
     min_{w,a} \quad 1/2(Dw - y)^TC^{-1}(Dw - y) +\gamma \sum_{i = 1}^p a_i \\
     s. t. \quad a_i \geq w_i, \quad a_i \geq -w_i
 \end{align}
where, $a$ is a new variable and $\gamma_{(p \times 1)}$ is the sparsity weight vector.\par To solve this optimization as a quadratic-programming using matlab, we should write it in the form of standard form:
\begin{align}
    min_x \quad x^THx + f^Tx \\
    s. t. \quad Ax \leq b
\end{align}
Therefore, in our problem, we have:
\begin{align}
 x_{(2p \times 1)}  = \left[ \begin{array}{c} w_{(p \times 1)}\\ a_{(p \times 1)} \end{array} \right] \\
 H_{(2p \times 2p)} = \left[ \begin{array}{cc} -(D^TC^{-1}D)_{(p \times p)} & 0\\0&0\end{array} \right] \\
 f_{(2p \times 1)} = \left[ \begin{array}{c}{(2y^TC^{-1}D)^T_{(p \times 1)}}\\\gamma_{(p \times 1)} \end{array} \right]\\
A_{(2p \times 2p)} = \left[ \begin{array}{cccccccccc}1&0&...&0&-1&0&...&0\\0&1&...&0&0&-1&...&0\\&&&&&&&&&\\0&0&...&1&0&0&...&-1\\
-1&0&...&0&-1&0&...&0\\0&-1&...&0&0&-1&...&0\\&&&&&&&&&\\0&0&...&-1&0&0&...&-1\\\end{array}\right]\\
b_{(2p \times 1)}=0_{(2p \times 1)}
\end{align}
%%%%%%%%%%%%%%%%%%%%%%%%%%%%%%%%%%%%%
%\subsubsection{Primary Results Without Adaptive Scheme}
%Fig.\ref{fig:primary_res} shows the estimated density obtained by the proposed method.
%\begin{figure}[hbt!]
%\centering
%\begin{tabular}{cc}
%\includegraphics[scale = 0.5]{Figures/primary_res_no_adaptive.PNG}     & \includegraphics[scale = 0.5]{Figures/gamma.PNG} \\
%     & 
%\end{tabular}
%\caption{primary results without adaptive setup}
%\label{fig:primary_res}
%\end{figure}
%%%%%%%%%%%%%%%%%%%%%%%%%%%%%%
\subsection{Considering Adaptive Scheme}
Taking into account the adaptive scheme means that we consider different bandwidths for each basis function $b_i(\cdot)$. Therefore, by considering $q$ fixed bandwidths for each basis function $b_i(\cdot)$, we have:\par%and adding a temporal laplacian regularization term~\cite{li2015temporal} to optimization problem, we have:
\begin{align}
     min_{w,a} \quad 1/2(Dw - y)^TC^{-1}(Dw - y) +\gamma \sum_{i = 1}^{pq} a_i \\% + \beta w^TLw \\
     s. t. \quad a_i \geq w_i, \quad a_i \geq -w_i\\
     %L_{ij}(pq \times pq)= \left\{ \begin{array}{cc} 1 & |i-j|<\frac{s}{2}\\ 0& o.w.\end{array} \right\\
    B_{(pq \times 1)} =[b_{11}, b_{21} ,..., b_{p1}, ..., b_{pq}]^T\\% \left[ \begin{array}{cccc}b_{11}& b_{21} &...& b_{p1}\\ b_{12}& b_{22} &...& b_{p2}\\&&&\\b_{1q}& b_{2q} &...& b_{pq}\end{array} \right]\\
    W_{(pq \times 1)} =[w_{11}, w_{21} ,..., w_{p1}, ..., w_{pq}]^T\\% \left[ \begin{array}{cccc}w_{11}& w_{21} &...& w_{p1}\\ w_{12}& w_{22} &...& w_{p2}\\&&&\\w_{1q}& w_{2q} &...& w_{pq}\end{array} \right]\\ 
    %\sum_{k=1}^qw_{ik}=1\\
    y_{(m \times 1)} = 1/N \sum_{i = 1}^N k(x_i)\\
    f_w(x) = \sum_{i=1}^{pq} w_{i}b_{i}(x) \\%= Tr((W^TB)_{(p \times p)})\\
    E(y) = \int k(x)f_w(x) dx= \left[\int k_{(m \times 1)}(x)b^T_{(1 \times pq)}(x)dx \right]w_{pq \times 1} \\
    =Dw\\
\end{align}
Therefore, we have:
\begin{align}
 min_x \quad x^THx + f^Tx \\
 s. t. \quad Ax \leq b\\
 x_{(2pq \times 1)}  = \left[ \begin{array}{c} w_{(pq \times 1)}\\ a_{(pq \times 1)} \end{array} \right] \\
 H_{(2pq \times 2pq)} = \left[ \begin{array}{cc} -(D^TC^{-1}D)_{(pq \times pq)} & 0\\0&0\end{array} \right] \\
 f_{(2pq \times 1)} = \left[ \begin{array}{c}{(2y^TC^{-1}D)^T_{(pq \times 1)}}\\\gamma_{(pq \times 1)} \end{array} \right]\\
A_{(2pq \times 2pq)} = \left[ \begin{array}{cccccccccc}1&0&...&0&-1&0&...&0\\0&1&...&0&0&-1&...&0\\&&&&&&&&&\\0&0&...&1&0&0&...&-1\\
-1&0&...&0&-1&0&...&0\\0&-1&...&0&0&-1&...&0\\&&&&&&&&&\\0&0&...&-1&0&0&...&-1\\\end{array}\right]\\
b_{(2pq \times 1)}=0_{(2pq \times 1)}
\end{align}
%%%%%%%%%%%%%%%%%%%%%%%%%%%%%%%%%%%%%
%\subsubsection{Results Obtained using Adaptive Scheme}
%Fig.\ref{fig:primary_res} shows the estimated density obtained by the proposed method.
%\begin{figure}[hbt!]
%\centering
%\begin{tabular}{cc}
%\includegraphics[scale = 0.5]{Figures/Adaptive/lin_apx_ml_s_2_beta_0.PNG}     & \includegraphics[scale = 0.5]{Figures/Adaptive/gamma_lin_apx_ml_s_2_beta_0.PNG} \\
 %    \multicolumn{2}{c}{$\beta = 0$} \\
%\includegraphics[scale = 0.7]{Figures/Adaptive/lin_apx_ml_s_2_beta_05.PNG}     & \includegraphics[scale = 0.5]{Figures/Adaptive/gamma_lin_apx_ml_s_2_beta_05.PNG} \\
 %    \multicolumn{2}{c}{$\beta = 0.5$} \\
%\includegraphics[scale = 0.6]{Figures/Adaptive/lin_apx_ml_s_2_beta_1.PNG} %    & \includegraphics[scale = 0.6]{Figures/Adaptive/gamma_lin_apx_ml_s_2_beta_1.PNG} \\
 %    \multicolumn{2}{c}{$\beta = 1$} 
%\end{tabular}
%\caption{primary results with adaptive scheme }
%\label{fig:primary_res}
%\end{figure}
%%%%%%%%%%%%%%%%%%%%%%%%%%%%%%
\section{SPADES~\cite{bunea2010spades} Method}
In this section we demonstrate the method proposed in~\cite{bunea2010spades}, which is a similar work to us. Suppose we have a set of $N$ sample points $X = [x_1, x_2, ..., x_N]$, independently drawn from a Gaussian mixture $f_w(x) = \sum_{i = 1}^{pq} b_i(x)w_i(x)$, which is a weighted combination of $p$ fixed Gaussian basis functions. For the first step, let's define the objective as to estimate the weights $w_{(pq \times 1)} = [w_1, w_2, ..., w_{pq}]^T$ in the $f_w$, knowing $pq$ fixed basis functions $b_{(pq\times 1)} = [b_1(\cdot), b_2(\cdot), ..., b_{pq}(\cdot)]^T$.\par, the objective in~\cite{bunea2010spades} is:
\begin{equation}
\begin{aligned}
    \min_{w, a} \left\{ -\frac{2}{N} \sum_{i=1}^N f_w(x_i) + ||f_w||^2 + \gamma\sum_{i=1}^{pq} a_i \right\}\\
    s. t. \quad a_i > w_i, a_i \geq -w_i
\end{aligned}   
\end{equation}
with
\begin{equation}
\begin{split}
||f_w(x)||^2 = \int (f_w(x))^2 dx = \int (w_jb_j(x))dx \\
= \int \sum_{j=1}^p \sum_{k=1}^p w_j w_kb_j(x) b_k(x)dx \\
= \sum_{j=1}^p \sum_{k=1}^p w_j w_k \int b_j(x) b_k(x) dx \\
= w^TCw
\end{split}
\end{equation}
and
\begin{equation}
c_{jk} = \int b_j(x) b_k(x) dx
\end{equation}
Using Matlab quadratic programming, we have:
\begin{equation}
    \begin{aligned}
    \min_x \quad x^THx + f^Tx\\
    s. t. \quad Ax \leq b \\
    x = \left[ \begin{array}{c} w_{pq \times 1}\\ a_{pq \times 1} \end{array} \right],\\
    H_{2pq \times 2pq} =  \left[ \begin{array}{cc} C_{pq \times pq} & 0\\0 & 0 \end{array}  \right], \\
    f_{2pq \times 1}=\left[ \begin{array}{c} -\frac{2}{N}\sum_{i=1}^N b_{pq \times 1}(x_i)\\\gamma_{pq \times 1}\end{array} \right], \\
    A_{(2pq \times 2pq)} = \left[ \begin{array}{cccccccccc}1&0&...&0&-1&0&...&0\\0&1&...&0&0&-1&...&0\\&&&&&&&&&\\0&0&...&1&0&0&...&-1\\
-1&0&...&0&-1&0&...&0\\0&-1&...&0&0&-1&...&0\\&&&&&&&&&\\0&0&...&-1&0&0&...&-1\\\end{array}\right]\\
b_{(2pq \times 1)}=0_{(2pq \times 1)}
    \end{aligned}
\end{equation}
